%-------------------------------------------------------
% Packages
%-------------------------------------------------------
\documentclass[a4paper, 11pt, twoside, numbers=noenddot, open=right, DIV=15, BCOR=1.3cm, listof=totocnumbered, bibliography=totocnumbered, appendixprefix=true, titlepage, headings=normal, parskip=relative]{scrbook}
\usepackage[ansinew]{inputenc}
\usepackage[T1]{fontenc}
\usepackage{mathptmx}
\linespread{1.1} 
\usepackage{floatflt} 
\usepackage{graphicx}
\usepackage[ngerman]{babel}
\usepackage{amsmath}
\usepackage{amssymb}
\usepackage[pdftex, bookmarks, bookmarksnumbered,pdfstartview=FitH]{hyperref}
\usepackage{colortbl}
\usepackage{float} 		
\usepackage{microtype}
\usepackage{sidecap}
\usepackage{cite}
\usepackage[numbers]{natbib}
\usepackage{picins}
\usepackage{dirtree}
\usepackage{wrapfig}
\usepackage{tabularx}
\usepackage{multirow}
\usepackage{multicol}
\usepackage{pdfpages}
\usepackage{enumitem}
\usepackage{trfsigns}
\usepackage{everysel}
\usepackage{subfig}
%\usepackage[plainpages=false, pdfpagelabels]{hyperref}

%-------------------------------------------------------
% Packete um Inkscape Schriften durch latex zu substituieren
%-------------------------------------------------------

\usepackage{import}
\graphicspath{{images/}} % Hier steht Pfad zu den Grafiken

%-------------------------------------------------------
% Darstellung �berschrift mit Serifen und Kapitel einrahmen
%-------------------------------------------------------
%\setkomafont{sectioning}{\bfseries}
\setkomafont{sectioning}{\slshape\bfseries}
%\renewcommand{\familydefault}{\sfdefault}
%\setkomafont{sectioning}{\bfseries}

%\renewcommand*{\chapterheadstartvskip}{\vspace*{-40pt}}

%\newcommand*{\ORIGchapterheadstartvskip}{}%
%\let\ORIGchapterheadstartvskip=\chapterheadstartvskip
%\renewcommand*{\chapterheadstartvskip}{%
%  \ORIGchapterheadstartvskip
%  {%
%    \setlength{\parskip}{0pt}%
%    \noindent\rule[-0.1\baselineskip]{\linewidth}{0.8pt}\par
%  }%
%}
\newcommand*{\ORIGchapterheadendvskip}{}%
\let\ORIGchapterheadendvskip=\chapterheadendvskip
\renewcommand*{\chapterheadendvskip}{%
  {%
    \setlength{\parskip}{0pt}%
    \noindent\rule[0.5\baselineskip]{\linewidth}{0.5pt}\par
  }%
  \ORIGchapterheadendvskip
}

%-------------------------------------------------------
% Matlab Code einbinden
%-------------------------------------------------------

\usepackage{listings}
\usepackage{color}
\usepackage[svgnames]{xcolor} 
\lstset{
   language = Matlab,
   basicstyle = {\ttfamily\scriptsize},
   keywordstyle = \color{black!80!black!100},
   identifierstyle = ,
   commentstyle = \color{black!50!black!100},
   stringstyle = \ttfamily,
   breaklines = true,
   numbers = left,
   numberstyle = \scriptsize,
   frame = single,
   backgroundcolor = \color{white!3}
}

%-------------------------------------------------------
% Zus�tzliches Layout Rand oben
%-------------------------------------------------------

\setlength{\topmargin}{-1.2 cm}      	  

%-------------------------------------------------------
%Absatzeinzug
%-------------------------------------------------------

\setparsizes{0pt}{4pt}{0pt plus 1fill}
%\setparsizes{0pt}{2pt}{0pt plus 1fill} 

%-------------------------------------------------------
% Kopf- und Fusszeile
%-------------------------------------------------------

\usepackage[automark]{scrpage2}
\pagestyle{scrheadings}

\clearscrheadfoot
\ohead{\headmark}
\ofoot[\pagemark]{\pagemark}
\setheadsepline{0.5pt}
%\setfootsepline{0.5pt}

%-------------------------------------------------------
% Tabellen Schattierung
%-------------------------------------------------------

\definecolor{Blue}{rgb}{0.0667,0.4000,0.6784}
\definecolor{Gray}{rgb}{0.9,0.9,0.9}

%-------------------------------------------------------
% Darstellung Formeln etc.
%-------------------------------------------------------

% F�r Matrizen und Vektoren
\newcommand{\m}[1]{\underline{#1}}
\newcommand{\mh}[1]{\underline{\hat{#1}}}
\newcommand{\mt}[1]{\underline{\tilde{#1}}}
\newcommand{\mb}[1]{\underline{\bar{#1}}}
\newcommand{\mhb}[1]{\underline{\hat{\bar{#1}}}}
\newcommand{\lt}[1]{\left\|{#1}\right\|_{2}}
\newcommand{\lm}[1]{\left\|{#1}\right\|_{\infty}}
\newcommand{\mr}[1]{#1}
\newcommand{\eq}[1]{\begin{align}{#1}\end{align}}
\newcommand{\seq}[1]{\begin{subequations}\begin{align}{#1}\end{align}\end{subequations}}

\definecolor{mainc}{RGB}{0,100,166} % Main Color ZHAW
\definecolor{subsc}{RGB}{206,0,60} % Subsidiary color School of Engineering
\definecolor{greyp}{RGB}{154,154,156} % Grey (pale)
\definecolor{greyd}{RGB}{74,93,104} % Alternate grey (dark)
\definecolor{grayd}{RGB}{100,100,100} % Alternate grey (dark)

% Color definitions for text and math mode for easy use
\newcommand{\tmc}[1]{\textcolor{mainc}{#1}}
\newcommand{\tsc}[1]{\textcolor{subsc}{#1}}
\newcommand{\tgd}[1]{\textcolor{greyd}{#1}}

\newcommand{\mmc}[1]{{\color{mainc} #1}}
\newcommand{\msc}[1]{{\color{subsc} #1}}
\newcommand{\mgd}[1]{{\color{greyd} #1}}



